\section{Discussion}
The main practical conclusion is that a mature fleet should be treated as a reusable source of controller-calibration information. In the studied LQI problem, a newly arriving target reaches near-oracle performance after roughly three local experiments when historical fleet observations are available, compared with more than six experiments under independent BO. The effect is especially pronounced in the tail: almost every warm-started target reaches the 5\% criterion within five evaluations, while fewer than half of independently optimized targets do so.

The more surprising result concerns heterogeneity. The vehicle families differ in mass, yaw inertia, and tire properties, and the identified dynamics are distinguishable. Nevertheless, those differences do not induce equally different rankings over the LQI weight space. Thus, \emph{plant heterogeneity need not imply optimization-task heterogeneity}. A representation useful for grouping dynamical models can therefore be unnecessary for transferring higher-level calibration experience.

This also explains why simple pooling can outperform more selective source use. In a three-dimensional search space with little target data, additional historical observations improve GP coverage even when they come from dynamically different clients. Soft dynamics weighting does not materially harm performance, but it provides no measurable advantage here. This should not be read as a universal rejection of dynamics-aware transfer; it establishes that the present fleet and LQI calibration objective are sufficiently transferable that explicit source selection is unnecessary.

Several limitations remain. Validation is simulation based; the nonlinear evaluation plant is intentionally mismatched to the controller-design model, but real-vehicle confirmation is missing. The study considers one controller structure, one three-dimensional calibration parameterization, and one lateral tracking application. Higher-dimensional or more task-specific calibration problems may exhibit less transferable landscapes. The architecture assumes scalar historical calibration records can be centrally stored or made available to the target and does not claim formal privacy. Finally, mature histories require prior fleet experimentation; their cost is amortized across future targets rather than eliminated.

\section{Conclusion}
This paper investigated whether controller-calibration experience accumulated by a mature heterogeneous fleet can reduce commissioning effort for a newly arriving client. In a pre-registered held-out study on nonlinear vehicle lateral dynamics, warm-starting Bayesian optimization with historical fleet observations reduced the mean number of target-local experiments required to reach within 5\% of the client-specific oracle from 6.63 to 3.07, while substantially reducing variability. Dynamics-aware source weighting did not improve over simple global pooling. A calibration-landscape analysis explains this result: despite substantial dynamics heterogeity, the ranking of LQI weight configurations remains highly consistent across clients. Mature fleet experience can therefore be broadly reusable for data-efficient controller calibration, and the heterogeneity relevant for plant modeling need not be the heterogeneity that matters for controller-parameter optimization.

\bibliographystyle{IEEEtran}
\bibliography{references}
\end{document}
